\chapter{Sesión 2: Backend Spring Boot base}
\label{ch:sesion2}

\section{Objetivo de la sesión}

Crear el proyecto backend en Spring Boot 4.1.1 con Java 21, definir las entidades JPA del dominio (\texttt{User}, \texttt{Channel}, \texttt{Post}, \texttt{Subscription}), sus repositorios, la migración Flyway inicial y la conexión a PostgreSQL, Redis y Kafka.

\section{Contexto}

La infraestructura ya estaba arriba (Sesión~\ref{ch:sesion1}). Faltaba el esqueleto del backend que consumiera esos servicios y validara el modelo relacional descrito en la Fase~2 de la metodología.

\section{Desarrollo}

\subsection{Generación del proyecto}

Se usó Spring Initializr vía \texttt{curl}. La primera tentativa con \texttt{bootVersion=3.3.5} falló porque esa versión ya no existe en el catálogo. El metadata disponible reportaba como versión por defecto \texttt{4.1.1.RELEASE}.

\begin{lstlisting}[language=bash]
cd ~/Escritorio/notiuam
curl -s https://start.spring.io/starter.tgz \
  -d type=maven-project \
  -d language=java \
  -d bootVersion=4.1.1.RELEASE \
  -d baseDir=backend \
  -d groupId=mx.uam \
  -d artifactId=notiuam-backend \
  -d name=notiuam-backend \
  -d packageName=mx.uam.notiuam \
  -d javaVersion=21 \
  -d dependencies=web,data-jpa,data-redis,security,validation,kafka,postgresql,lombok,flyway \
  -o backend.tgz
tar -xzvf backend.tgz
rm backend.tgz
\end{lstlisting}

\begin{aviso}[Nota sobre alias]
En Spring Boot~4 el starter web se llama \texttt{spring-boot-starter-webmvc}, pero el \textbf{alias del Initializr sigue siendo \texttt{web}}. El alias \texttt{webmvc} no es válido y produce error \texttt{Invalid dependency}.
\end{aviso}

\subsection{Dependencias adicionales}

Se agregaron a \texttt{pom.xml} las librerías para manejo de JWT y para documentación OpenAPI:

\begin{lstlisting}[language=xml]
<!-- JWT -->
<dependency>
  <groupId>io.jsonwebtoken</groupId>
  <artifactId>jjwt-api</artifactId>
  <version>0.12.6</version>
</dependency>
<dependency>
  <groupId>io.jsonwebtoken</groupId>
  <artifactId>jjwt-impl</artifactId>
  <version>0.12.6</version>
  <scope>runtime</scope>
</dependency>
<dependency>
  <groupId>io.jsonwebtoken</groupId>
  <artifactId>jjwt-jackson</artifactId>
  <version>0.12.6</version>
  <scope>runtime</scope>
</dependency>

<!-- OpenAPI / Swagger UI -->
<dependency>
  <groupId>org.springdoc</groupId>
  <artifactId>springdoc-openapi-starter-webmvc-ui</artifactId>
  <version>2.6.0</version>
</dependency>
\end{lstlisting}

\subsection{Entidades JPA}

Se modelaron cuatro entidades correspondientes a las tablas del protocolo: \texttt{users}, \texttt{channels}, \texttt{posts} y \texttt{subscriptions}.

\subsubsection{\texttt{User}}

\begin{lstlisting}[language=java]
@Entity
@Table(name = "users",
       uniqueConstraints = @UniqueConstraint(columnNames = "email"))
@EntityListeners(AuditingEntityListener.class)
@Getter @Setter @NoArgsConstructor @AllArgsConstructor @Builder
public class User {

    @Id @GeneratedValue(strategy = GenerationType.UUID)
    private UUID id;

    @Column(nullable = false, length = 180)
    private String email;

    @Column(name = "password_hash", nullable = false, length = 100)
    private String passwordHash;

    @Column(name = "display_name", nullable = false, length = 120)
    private String displayName;

    @Enumerated(EnumType.STRING)
    @Column(nullable = false, length = 20)
    private Role role;

    @Column(nullable = false)
    @Builder.Default
    private boolean enabled = true;

    @CreatedDate @Column(name = "created_at", nullable = false, updatable = false)
    private Instant createdAt;

    @LastModifiedDate @Column(name = "updated_at")
    private Instant updatedAt;
}
\end{lstlisting}

\subsubsection{Enumeraciones del dominio}

\begin{lstlisting}[language=java]
public enum Role { STUDENT, PROFESSOR, ADMIN }

public enum ChannelCategory {
    ACADEMIC, CULTURAL, SPORTS, ADMINISTRATIVE, SOCIAL, OTHER
}
\end{lstlisting}

\subsubsection{\texttt{Channel}, \texttt{Post} y \texttt{Subscription}}

Las tres entidades siguen el mismo patrón: UUID como clave primaria, relaciones \texttt{@ManyToOne} hacia \texttt{User} y \texttt{Channel} cuando aplica, y campos de auditoría con \texttt{@CreatedDate}.

La entidad \texttt{Subscription} usa una \textbf{clave primaria compuesta} (\texttt{userId}, \texttt{channelId}) implementada con \texttt{@IdClass} y una clase auxiliar \texttt{SubscriptionId}.

\subsection{Repositorios Spring Data}

Cuatro interfaces que extienden \texttt{JpaRepository}:

\begin{lstlisting}[language=java]
public interface UserRepository extends JpaRepository<User, UUID> {
    Optional<User> findByEmail(String email);
    boolean existsByEmail(String email);
}

public interface ChannelRepository extends JpaRepository<Channel, UUID> {
    List<Channel> findByPublicChannelTrueOrderByCreatedAtDesc();
    List<Channel> findByCategoryAndPublicChannelTrue(ChannelCategory category);
    List<Channel> findByNameContainingIgnoreCaseAndPublicChannelTrue(String q);
    List<Channel> findByOwnerId(UUID ownerId);
}

public interface PostRepository extends JpaRepository<Post, UUID> {
    Page<Post> findByChannelIdOrderByCreatedAtDesc(UUID channelId, Pageable pageable);
}

public interface SubscriptionRepository
        extends JpaRepository<Subscription, SubscriptionId> {
    List<Subscription> findByUserId(UUID userId);
    List<Subscription> findByChannelId(UUID channelId);
    boolean existsByUserIdAndChannelId(UUID userId, UUID channelId);
    long countByChannelId(UUID channelId);
}
\end{lstlisting}

\subsection{Migración Flyway}

Se eligió Flyway en lugar de \texttt{ddl-auto=update} porque el protocolo pide trazabilidad del esquema (Fase~2). El archivo \texttt{V1\_\_init.sql} crea las cuatro tablas con sus índices:

\begin{lstlisting}[language=SQL]
CREATE EXTENSION IF NOT EXISTS "pgcrypto";

CREATE TABLE users (
    id              UUID PRIMARY KEY DEFAULT gen_random_uuid(),
    email           VARCHAR(180) NOT NULL UNIQUE,
    password_hash   VARCHAR(100) NOT NULL,
    display_name    VARCHAR(120) NOT NULL,
    role            VARCHAR(20)  NOT NULL,
    enabled         BOOLEAN      NOT NULL DEFAULT TRUE,
    created_at      TIMESTAMPTZ  NOT NULL DEFAULT NOW(),
    updated_at      TIMESTAMPTZ
);

CREATE TABLE channels (
    id           UUID PRIMARY KEY DEFAULT gen_random_uuid(),
    name         VARCHAR(120) NOT NULL,
    description  VARCHAR(500),
    category     VARCHAR(30)  NOT NULL,
    owner_id     UUID         NOT NULL REFERENCES users(id) ON DELETE CASCADE,
    is_public    BOOLEAN      NOT NULL DEFAULT TRUE,
    created_at   TIMESTAMPTZ  NOT NULL DEFAULT NOW()
);

CREATE INDEX idx_channels_public    ON channels(is_public);
CREATE INDEX idx_channels_category  ON channels(category);
CREATE INDEX idx_channels_owner     ON channels(owner_id);

CREATE TABLE posts (
    id          UUID PRIMARY KEY DEFAULT gen_random_uuid(),
    channel_id  UUID          NOT NULL REFERENCES channels(id) ON DELETE CASCADE,
    author_id   UUID          NOT NULL REFERENCES users(id)    ON DELETE CASCADE,
    content     VARCHAR(2000) NOT NULL,
    created_at  TIMESTAMPTZ   NOT NULL DEFAULT NOW()
);

CREATE INDEX idx_posts_channel_created ON posts(channel_id, created_at DESC);

CREATE TABLE subscriptions (
    user_id        UUID        NOT NULL REFERENCES users(id)    ON DELETE CASCADE,
    channel_id     UUID        NOT NULL REFERENCES channels(id) ON DELETE CASCADE,
    subscribed_at  TIMESTAMPTZ NOT NULL DEFAULT NOW(),
    PRIMARY KEY (user_id, channel_id)
);

CREATE INDEX idx_subs_channel ON subscriptions(channel_id);
\end{lstlisting}

\subsection{Configuración (\texttt{application.yml})}

Se conectó el backend a los tres servicios de la Sesión~1:

\begin{lstlisting}[language=yaml]
spring:
  datasource:
    url: jdbc:postgresql://localhost:5432/notiuam
    username: notiuam
    password: notiuam_dev_pass
    hikari:
      maximum-pool-size: 10
      minimum-idle: 2

  jpa:
    hibernate:
      ddl-auto: validate          # Flyway maneja el esquema
    open-in-view: false

  flyway:
    enabled: true
    locations: classpath:db/migration

  data:
    redis:
      host: localhost
      port: 6379
      repositories:
        enabled: false            # Evita el warning "Could not safely identify store"

  kafka:
    bootstrap-servers: localhost:9092
    producer:
      key-serializer: org.apache.kafka.common.serialization.StringSerializer
      value-serializer: org.springframework.kafka.support.serializer.JsonSerializer
    consumer:
      group-id: notiuam
      auto-offset-reset: earliest

server:
  port: 8081

jwt:
  secret: ${JWT_SECRET:change_me_in_prod}
  expiration-ms: 3600000
  refresh-expiration-ms: 604800000
\end{lstlisting}

\subsection{Clase principal}

\begin{lstlisting}[language=java]
@SpringBootApplication
@EnableJpaAuditing
public class NotiuamApplication {
    public static void main(String[] args) {
        SpringApplication.run(NotiuamApplication.class, args);
    }
}
\end{lstlisting}

\section{Verificación}

Arrancar la aplicación:

\begin{lstlisting}[language=bash]
cd ~/Escritorio/notiuam/backend
./mvnw spring-boot:run
\end{lstlisting}

En la salida deben verse, en este orden:

\begin{enumerate}[leftmargin=*]
  \item \texttt{HikariPool-1 - Start completed} (conexión a PostgreSQL).
  \item \texttt{Successfully applied 1 migration to schema "public"} (Flyway).
  \item \texttt{Initialized JPA EntityManagerFactory} (Hibernate validó las entidades).
  \item \texttt{Tomcat started on port 8081}.
  \item \texttt{Started NotiuamApplication in X seconds}.
\end{enumerate}

Verificar las tablas:

\begin{lstlisting}[language=bash]
docker exec -it notiuam-postgres psql -U notiuam -d notiuam -c "\dt"
\end{lstlisting}

\begin{exito}[Resultado esperado]
Las tablas \texttt{users}, \texttt{channels}, \texttt{posts}, \texttt{subscriptions} y \texttt{flyway\_schema\_history}.
\end{exito}

\section{Resultado}

\begin{exito}[Estado al cierre]
\begin{itemize}[leftmargin=*]
  \item Backend compila y arranca en el puerto 8081.
  \item Modelo relacional creado y validado por Hibernate.
  \item Conexión a PostgreSQL, Redis y Kafka configurada.
  \item Aún sin endpoints ni seguridad: Spring Security generaba una contraseña temporal de desarrollo.
\end{itemize}
\end{exito}

\textbf{Pendiente para la siguiente sesión:} implementar autenticación JWT con endpoints de registro, login y refresh.